\chapter{Chercher dans la littérature sans s'attribuer de priorité}\label{ch:litterature}

Un assistant lit vite et cherche loin. Il est tentant de lui demander si un énoncé est
nouveau. Le danger n'est pas une formule fausse : c'est une affirmation fausse sur le
travail des autres, du type « personne n'avait démontré ceci ».

\section{Une affirmation sur la littérature, pas sur l'énoncé}

\begin{regle}
« Nouveau » n'est jamais une conclusion. On écrit « non trouvé dans les sources
suivantes », et l'on nomme les sources, avec la section.
\end{regle}

On peut lire un énoncé ; on ne peut que \emph{chercher} dans la littérature, jamais
l'épuiser. Une liste nommée de ce qui a été cherché se réfute : un lecteur peut aller voir.
« Pas dans la littérature standard » ne se réfute pas, et ne vaut donc rien.

\section{Des statuts fermés}

Chaque énoncé candidat reçoit un statut pris dans une liste courte :
\begin{description}
  \item[non cherché] personne n'a regardé ; c'est le statut honnête par défaut ;
  \item[candidat] cherché dans les sources nommées, non trouvé ;
  \item[trouvé] présent dans la littérature, avec la référence ; on garde la ligne, pour
        épargner la recherche au suivant ;
  \item[réfuté] faux tel qu'énoncé, par un contre-exemple ; on garde la ligne, avec ce
        qui tient encore ;
  \item[confirmé] vérifié par une personne. Seule une personne pose ce statut.
\end{description}

Une ligne « trouvé » est un bon résultat : elle montre que la liste est élaguée au lieu de
grossir.

\section{Séparer le texte et nos ajouts}

Une lecture moderne complète souvent l'original : une hypothèse sous-entendue, un pas
sauté. Si l'énoncé intéressant n'existe qu'avec ce complément, il est en partie de nous, et
l'entrée le dit. Un énoncé qui repose sur un mot illisible ne peut pas être un candidat.

\section{Vérifier les citations}

Un modèle de langage invente des références plausibles : bon auteur, bonne revue, titre
faux, ou théorème mal numéroté. Chaque référence se vérifie sur une source canonique
(l'éditeur, arXiv, une base bibliographique), et chaque affirmation qu'on attribue à un
article se vérifie dans l'article, à la page.

\begin{exemple}
Un brouillon attribuait à un article la démonstration d'un principe général. La lecture des
deux premières pages a montré que les auteurs parlaient d'une « philosophie » et
démontraient des énoncés plus restreints. Les trois documents du projet qui reprenaient
l'affirmation ont été corrigés dans la foulée.
\end{exemple}
